%%
%% This is file `fillwith-example.tex',
%% generated with the docstrip utility.
%%
%% The original source files were:
%%
%% fillwith.dtx  (with options: `ee')
%% Copyright (C) 2024 Clea F. Rees.
%% 
%% This work may be distributed and/or modified under the
%% conditions of the LaTeX Project Public License, either version 1.3c
%% of this license or (at your option) any later version.
%% The latest version of this license is in
%%   https://www.latex-project.org/lppl.txt
%% and version 1.3c or later is part of all distributions of LaTeX
%% version 2008-05-04 or later.
%% 
%% This work has the LPPL maintenance status `maintained'.
%% 
%% The Current Maintainer of this work is Clea F. Rees.
%% 
%% This file may only be distributed together with a copy of the package
%% fillwith. You may however distribute the package fillwith without
%% such generated files.
%% 
%% This work consists of all files listed in manifest.txt.
%%%%%%%%%%%%%%%%%%%%%%%%%%%%%%%%%%%%%%%%%%%%%%%%%
%%%%%%%%%%%%%%%%%%%%%%%%%%%%%%%%%%%%%%%%%%%%%%%%%
\listfiles
\documentclass[british,a6paper]{article}
\usepackage{babel}
\usepackage[scale=.85,marginparwidth=0pt,marginparsep=0pt]{geometry}
\usepackage{fillwith}
\usepackage{cfr-lm}
\begin{document}
\author{A. N. Author}
\title{\texttt{fillwith} Demo}
\maketitle
\begin{enumerate}
  \item What is the answer to the question?
    \fillwithrules
  \item How did you answer the previous question?
    \fillwith*[style=line]
  \item Why answer?

    \fillwith*[style=rule,fillwith ht=1.25,colour=blue!40!cyan]
    \clearpage
  \item Evaluate your responses so far.
    \fillwithdottedlines
  \item Evaluate your responses so far.

    {%
      \fillwithset{dotted colour=magenta}%
      \fillwithdottedlines[1.5]
    }
  \item Criticise your evaluation.
    \fillwithdottedlines*
    \clearpage
  \item Assess your answers.
    \fillwithnolines[goal ht = .33\textheight]
  \item Assess your answers.

    \fillwithnolines[style=rules,no font=\bfseries\tiny,colour=green!50!gray,goal ht = .67\textheight]
  \item Summarise your analysis.
    \fillwithnolines*
\end{enumerate}

\end{document}
\endinput
%%
%% End of file `fillwith-example.tex'.
